% ================================================================
% SECTION 7: DISCRETIZATION MODULE
% ================================================================
\section{Discretization Module}
\label{sec:discretization}

\begin{center}
\code{bionetflux.core.discretization}
\end{center}

This module (133 lines) provides classes for spatial mesh generation and time discretization management.
It does \emph{not} store matrices or polynomial bases---those live in
\code{ElementaryMatrices} (Section~\ref{sec:elementary_matrices}) and the static-condensation classes (Section~\ref{sec:static_condensation}).

% ------------------------------------------------------------------
\subsection{compute\_n\_elements\_from\_h}

\begin{lstlisting}[language=Python, caption={Utility: mesh size to element count}]
def compute_n_elements_from_h(domain_length: float, h: float) -> int:
    """
    Return the *even* integer closest to domain_length / h,
    with a minimum of 4.
    Being even guarantees the segment midpoint is a mesh node.
    """
\end{lstlisting}

\noindent Raises \code{ValueError} if \code{domain\_length \textless= 0} or \code{h \textless= 0}.

% ------------------------------------------------------------------
\subsection{Discretization (single domain)}

\begin{lstlisting}[language=Python, caption={Discretization constructor}]
class Discretization:
    def __init__(self, n_elements: int, domain_start: float = 0.0,
                 domain_length: float = 1.0, stab_constant: float = 1.0):
\end{lstlisting}

\noindent Computed at construction (via \code{\_generate\_mesh}):

\begin{longtable}{p{5cm}p{9.5cm}}
\toprule
\textbf{Attribute} & \textbf{Description} \\
\midrule
\code{n\_nodes} & \code{n\_elements + 1} \\
\code{element\_length} & \code{domain\_length / n\_elements} \\
\code{nodes} & \code{np.linspace(domain\_start, domain\_start + domain\_length, n\_nodes)} \\
\code{elements} & Integer array of shape (\code{n\_elements}, 2), each row is \code{[i, i+1]} \\
\code{element\_centers} & Midpoints, \code{nodes[:-1] + element\_length / 2} \\
\bottomrule
\end{longtable}

\noindent Additional methods:

\begin{longtable}{p{6.5cm}p{8cm}}
\toprule
\textbf{Method} & \textbf{Description} \\
\midrule
\code{get\_mesh\_info() -> dict} & Returns a dictionary with all mesh parameters \\
\code{set\_tau(tau\_values: List[float])} & Set per-equation stabilization parameters; stored in \code{self.tau} \\
\bottomrule
\end{longtable}

% ------------------------------------------------------------------
\subsection{GlobalDiscretization (multi-domain + time)}
\label{sec:global_disc}

\begin{lstlisting}[language=Python, caption={GlobalDiscretization constructor}]
class GlobalDiscretization:
    def __init__(self, spatial_discretizations: List[Discretization]):
\end{lstlisting}

\noindent Computed at construction:
\begin{itemize}
  \item \code{n\_domains}: number of domains.
  \item \code{total\_elements}, \code{total\_nodes}: sums across all domains.
\end{itemize}

\noindent Time parameters (initially \code{None}):

\begin{longtable}{p{6.5cm}p{8cm}}
\toprule
\textbf{Method / Attribute} & \textbf{Description} \\
\midrule
\code{set\_time\_parameters(dt, T)} & Compute \code{n\_time\_steps = ceil(T/dt)} and fill \code{time\_points = linspace(0, T, n\_time\_steps+1)} \\
\code{get\_spatial\_discretization(i)} & Return the $i$-th \code{Discretization} \\
\code{get\_time\_info() -> dict} & Dict with \code{dt}, \code{T}, \code{n\_time\_steps}, \code{time\_points} \\
\code{get\_global\_info() -> dict} & Dict with \code{n\_domains}, \code{total\_elements} \\
\bottomrule
\end{longtable}
